\sectionguard
\section{Scope, Status, and Qualifications}

\subsection*{Question and thesis boundary}

This study answers one question: what was reported at La Salette on
19 September 1846 on the nearest record this project could reach;
how the public discourse is transmitted and where its strata differ;
what exactly the Bishop of Grenoble judged on 19 September 1851, in
which words and with what object; what happened to the
secrets---their writing, their delivery to Pius~IX, their expansion in
publication, and the Roman acts of 1880, 1915, 1916, and 1923;
whether the text published in 1879 is the text sent to Pius~IX in
1851; how the
shrine, the Missionaries, the liturgy, and the popes received the
event; and what all of that obliges and leaves open. It is a
source-audited study instrument, ecclesial-judgment-first in method:
not a devotional narrative, a critical edition, a magisterial act, a
biography, or a complete history of La Salette. It proposes no
judgment of its own on supernaturality; the judgment it does propose,
at \S6.9, is about the textual identity of one nineteenth-century
document and nothing else.

\subsection*{The secret: what is published, and on what disclosure}

This edition publishes and analyzes the text printed at Lecce in
1879 and reprinted at Louvain in 1881, and reaches a judgment on its
relation to the manuscripts written for Pius~IX in July 1851. The
governing disclosure is at \S5.1 and is not repeated here: the Holy
Office's decree of 21 December 1915 commanded all the faithful of
every region not to discuss or treat of the subject, with penalties
and with an express saving clause for the devotion; the acts of 1916
and 1923 applied it to named publications; the decree's present
force is genuinely unresolved, this project's bounded searches
located no act abrogating or modifying it, and this edition asserts
neither that it has lapsed nor that it binds. The analysis is
published on the express direction of the project's maintainer,
recorded as the editorial act of a private study. Nothing in \S5 or
\S6 is an ecclesiastical interpretation, an authentic interpretation
of any canon or decree, a dispensation, a canonical opinion, or
advice to any other author. The judgment reached at \S6.9---that the
1879 text is not the 1851 text---is project synthesis, is labelled
as such at its own locus, has had no ecclesiastical or independent
scholarly review, and is stated there with its confidence levels,
its strongest counterargument, and the evidence that would change
it. Throughout, the \emph{public message} of 1846, which is approved
in its event and touched by no Roman restriction, is kept distinct
from the \emph{secrets}, which are approved by nobody.

\subsection*{Corpus and exclusions}

Included: the single reported apparition of 19 September 1846 to
Maximin Giraud and Mélanie Calvat; the transmission strata named at
\S2.1; the public discourse in its 1854-printed recital form and in
the received custodial form, compared element by element; the 1847
depositions and the commission as the 1851 act describes them; the
mandement of 19 September 1851 at its operative clause and article
structure; the mandement of 1 May 1852 founding the Missionaries;
the Roman grants of 1852; the Rites rescripts of 12 May and 7 July
1877; the documentary chronology of the secrets, the reported letter
of the Secretary of the Holy Office of 14 August 1880, and the Holy
Office and Index acts of 1915, 1916, and 1923; the text of the 1879
Lecce booklet as the 1881 Louvain reprint prints it, published at
its principal elements, analyzed, and compared with the 1846
discourse; the seers' later lives at
the custodian's and the 1854 witness's ceilings, including the Ars
episode; Ginoulhiac's acts of 1854--1857 and his 1855 sentence; papal
reception in 1945, 1961, and 1996; the 2016 inscription in the
Proper of France; and current standing under the 2024 norms.

Excluded or bounded: the whole text of the 1879 booklet, which is
not reproduced integrally---its elements are inventoried and its
load-bearing sentences quoted; any text of Maximin Giraud's secret,
none having been reached at an identified imprint; the wording of
the 1880 letter, not reached in any language; any judgment
on supernaturality; any certification of a cure, prodigy, or
property of the spring; any decoding, dating, or application of the
dearth prophecy or of the saying about children under seven; the
identification hypothesis about the figure seen, whose litigation is
recorded as existing and not adjudicated; the personal sanctity,
culpability, mental state, or cause of death of either seer; and
every claim of the ``all approved apparitions'' type---this is a
single-event study and asserts no registry.

\subsection*{Geography, chronology, and currentness}

Subject geography: the mountain above La Salette and the village of
Corps in the Isère; Grenoble; Rome; Lecce and Castellammare di
Stabia; Louvain; Ars and Lyons. Subject range: 1846 to the acts
current at the as-of date. As-of date: all mutable claims---current
controlling status, the Dicastery's document index, the shrine and
institute pages, the diocesan ordinary, the calendar entry---were
checked on 2026-07-25. Stable historical facts are not given
artificial currentness. The 1852 indulgence and Mass privileges are
reported as acts of their own period and are expressly \emph{not}
stated as discipline in force. Negative results are bounded searches,
correctable by better evidence, never guarantees.

\subsection*{Taxonomy and terminology}

Kept distinct throughout: \emph{recital stratum} (the children's
narrative as the diocesan documents transmit it, read here in the
1854 printing); \emph{received custodial text} (the harmonized
French discourse published by the shrine and the institute);
\emph{period apologetic literature} (the 1852 and 1854 English
volumes, the 1866 Grenoble biography); \emph{hostile literature}
(the 1855 volume, used only for the existence of the litigation);
\emph{promoter's testimony} (Bishop Zola's 1880 letter and the 1882
Roubaud pamphlet, used for the custody chronology and for their
authors' admissions against interest, never for their evaluations);
\emph{analyzed text} (the 1879 Lecce booklet as the 1881 Louvain
reprint prints it: published, quoted, and assessed, and asserted by
this study neither as true nor as revealed);
\emph{competent judgment} (the 1851
mandement); \emph{disciplinary act} (the 1854 mandements; the 1915,
1916, and 1923 Roman acts); \emph{act of cult} (the 1852 grants; the
1877 rescripts; the 2016 memorial); \emph{papal reception} (1945,
1961, 1996); and \emph{editorial synthesis}, always labelled.
``Reported'' is the default register for every word and phenomenon
attributed to the apparition. \term{Fondés à la croire indubitable
et certaine} is the 1851 act's own formula and is not translated into
the 2024 taxonomy. \term{Fides humana}---in the 1877 rescript's own
words, \term{fides solum humana}---names the assent invited by
permitted private revelation.

\subsection*{Source hierarchy and method}

Priority: the exact operative texts in identified official witnesses
(\work{AAS} 1 (1909), 7 (1915), 8 (1916), 15 (1923), 18 (1926), 38
(1946), 53 (1961); \work{ASS} 11 (1878), 13 (1880); the Holy See's
web texts for 1996 and 2024), each read directly in the Holy See's
own archive on 2026-07-25; then the custodial institutions (the
Sanctuary, the Missionaries, the French episcopal conference's
official site) for what they alone witness, at
institutional-self-description ceilings; then identified
nineteenth-century printings for the mandement, the discourse, and
the documentary chronology, each named with its imprint and its
partisanship stated. Quotations are exact from the named witness;
French, Latin, and Italian are quoted in the original where wording
is operative, with editorial translations visibly editorial. No
project or AI rendering is offered as a received translation, and no
text is printed for recitation.

\subsection*{Qualifications and unresolved items}

(1)~\textbf{No French printing of the integral 1851 mandement was
located.} Article~1 is therefore quoted in French only as far as the
Sanctuary's transcription (which carries its own ellipsis) and the
Dicastery's 2024 citation print it; the whole article is given in an
1854 English rendering and the eight-article structure from an 1854
Italian translation of the pilgrim's manual. No French reconstruction
is offered anywhere in this study. (2)~The 1852 English translation
of the mandement's preamble breaks off before the operative
articles; the two witnesses for the articles differ in fullness, and
the article inventory follows the Italian witness with the difference
recorded. (3)~The critical documentary edition of the early La
Salette records, the diocesan archives of Grenoble, the
\term{rapporteur}'s report and its supplement, Maximin's pamphlet,
and Mélanie's autobiography were not consulted. (4)~The two 1851
manuscripts were not reached in any form; later reports of their
recovery and facsimile publication were not verified and nothing
rests on them; the analysis of \S6 therefore compares the 1879 text
with the \emph{1846} record, which exists, and not with the 1851
text, which does not. (5)~No copy or image of the 1879 Lecce imprint
was reached. Everything published and analyzed in \S6 is read in the
1881 Louvain reprint, whose whole text was collated against that
copy's page images on 2026-07-25; the relation of the 1881 setting
to the 1879 setting rests on the reprint's own claim and on Bishop
Zola's letter and was not independently verified. A \emph{correction
to an earlier state of this record}: the 15\slash 18 November 1879
discrepancy formerly preserved between the printed formula and
Zola's letter does not exist. Page-image collation shows both at 15
November 1879; the ``18'' was a defect of the scan's text layer.
Two smaller items in the same volume remain open and are not
resolved: an ambiguous two-digit number at p.~9 (33 or 35), and the
title page's ``depuis 1867'' for Zola's direction against his own
``de 1868 jusqu'en 1873.'' (5a)~No printing, transcription, or
facsimile of any text of Maximin Giraud's secret was reached at an
identified imprint; none is published, quoted, or characterized.
(5b)~The letter of the Secretary of the Holy Office of 14 August
1880 was not reached in any language; its existence, date, author,
and substance rest on one partisan witness of 1882 who states that
he saw only a translation, and no clause of this study depends on
its wording. (5c)~The \work{Traité de la vraie dévotion} of Louis
Marie Grignion de Montfort, whose vocabulary \S6.5 notices in the
1879 text, was not consulted in any edition; the observation is a
correspondence of vocabulary and no literary dependence is asserted.
(6)~The Holy See's \work{AAS} and
\work{ASS} archive files are typeset text transcriptions without
page images, so no image collation of the 1915 or 1923 decrees was
possible; two evident transcription defects are bracketed or
recorded (\term{cum} for \term{tum}; \term{Beconciliatricis} for
\term{Reconciliatricis}; the pamphlet year printed ``1845''), and a
third in an 1847 deposition (``1840'' for 1846). (7)~The \work{AAS}
sweep covered volumes 1 (1909) through 98 (2006) except volume 75
(1983), which that route does not serve; volumes from 99 (2007) are
not published there as annual files and were not searched, and the
1966 Index notification's locus is taken from an existing repository
source record rather than re-verified. (8)~The 2016 decree
inscribing the memorial in the Proper of France was not located as a
document; its protocol number, Latin text, and rank terminology are
unverified and the act is reported at the French episcopal
conference's institutional ceiling. (9)~The Sanctuary's biographical
page for the seers is cited from an Internet Archive capture of 8
April 2013 because the current site no longer carries it; the
Sanctuary's own datings, including Mélanie's death date, are reported
as that page gives them and no civil or parish record was consulted.
(10)~The 1855 Ginoulhiac sentence is cited as the Sanctuary's
received formula; no printed 1855 text was reached. (11)~The Ars
episode is reported from a pro-La Salette 1854 witness printing
Maximin's own reported answers; no testimony of Jean-Marie Vianney
was reached, and nothing said in a confessional is adjudicated.
(12)~No membership statistics for the Missionaries are asserted;
none was found stated with a date on their own pages. (12a)~The 1882
Roubaud pamphlet is an openly partisan defence of the secret, used
only for the facts it concedes against its own case---the Index's
refusal to condemn, the Holy Office's taking up of the matter, the
letter of 14 August 1880, and the state of the 1882
controversy---and never for its arguments or evaluations; the
periodicals it names were not consulted, and the solemn coronation
of 20 August 1879 that it and the 1881 \term{avant-propos} report
was not located as an act. (13)~No
independent Mariological, historical, canonical, archival, French or
Occitan philological, liturgical, or textual review has occurred; no
ecclesiastical review or approval is claimed or implied.

\subsection*{Rights}

Project prose is original and intended for CC~BY~4.0. Quotations
remain under their own status. The 1851 and 1852 mandements, the
nineteenth-century French, Latin, and Italian printings, and the
Holy See acts published in \work{Acta Sanctae Sedis} and
\work{Acta Apostolicae Sedis} through 1961 are public domain in
text and are quoted with citation; the Holy See's web texts of the
1996 letters and the 2024 norms are \copyright\ Libreria Editrice
Vaticana and are quoted within ordinary scholarly limits with
citation; the Sanctuary of Notre-Dame de La Salette, the
Missionaries of Our Lady of La Salette, the Diocese of
Grenoble--Vienne, and the French episcopal conference's Nominis
service are \copyright\ their publishers and are quoted briefly with
citation. The 1881 Louvain imprint and the 1882 Roubaud pamphlet are
public domain by age; the French quoted from them in \S5 and \S6 is
their own wording, and the English renderings beside it are
editorial commentary of this project, not received translations.
Neither volume is reproduced integrally. The received custodial text
of the discourse is likewise not reproduced in full: only its opening
and closing lines are quoted and the custodial publishers are cited
for the whole. Online
availability is not treated as a reuse license; no complete document
is reproduced; scans consulted for reading are not redistributed.
